\chapter{Voorrangsregels}\label{ch:g7:priorities}

Hoeveel is $3 + 4 \times 5$? Reken je van links naar rechts, dan krijg je $35$;
vermenigvuldig je eerst, dan krijg je $23$. De wiskunde heeft één enkel antwoord
nodig, en dus zijn er verkeersregels voor berekeningen --- de voorrangsregels.
Dit hoofdstuk zet ze op een rij en introduceert een eerste algebraïsche
eigenschap: de distributiviteit.

\section{De voorrangsregels}

\begin{theorem}[Voorrangsregels]\label{thm:g7:priorities:rules}
In een berekening zonder haakjes:
\begin{enumerate}
  \item worden vermenigvuldigingen en delingen \emph{voor} optellingen en
        aftrekkingen uitgevoerd;
  \item worden bewerkingen van hetzelfde niveau van links naar rechts
        uitgevoerd.
\end{enumerate}
Haakjes gaan boven alles: wat tussen haakjes staat, wordt eerst uitgerekend, te
beginnen bij de binnenste.
\end{theorem}

\begin{example}\label{ex:g7:priorities:basic}
Eén stap per regel, met (in gedachten) een streep onder wat je uitrekent:
\begin{align*}
3 + 4 \times 5 &= 3 + 20 = 23
&&\text{(eerst de vermenigvuldiging)}\\
20 - 8 + 2 &= 12 + 2 = 14
&&\text{(zelfde niveau: van links naar rechts)}\\
18 \div 3 \times 2 &= 6 \times 2 = 12
&&\text{(zelfde niveau: van links naar rechts)}\\
(3 + 4) \times 5 &= 7 \times 5 = 35
&&\text{(eerst de haakjes).}
\end{align*}
Let op de valstrikken: $20 - 8 + 2$ is \emph{niet} $20 - 10$, en
$18 \div 3 \times 2$ is \emph{niet} $18 \div 6$.
\end{example}

\begin{example}[Haakjes in haakjes]\label{ex:g7:priorities:nested}
\begin{align*}
50 - 2 \times \bigl(3 + (10 - 4)\bigr)
&= 50 - 2 \times (3 + 6) && \text{(binnenste haakje)}\\
&= 50 - 2 \times 9 && \text{(resterend haakje)}\\
&= 50 - 18 && \text{(vermenigvuldiging)}\\
&= 32 .
\end{align*}
\end{example}

\begin{remark}[In de breukstreep zitten haakjes verstopt]
Een breukstreep groepeert haar \omterm{def:g4:fractions:def}{teller} en haar \omterm{def:g4:fractions:def}{noemer}, alsof elk van beide
onzichtbare haakjes heeft:
\[
\frac{12 + 8}{4} = \frac{20}{4} = 5,
\qquad\text{en niet}\quad 12 + \frac{8}{4} = 14 .
\]
\end{remark}

\begin{method}[Veilig rekenen]\label{met:g7:priorities:safe}
\begin{enumerate}
  \item Zoek de bewerking die als eerste moet (binnenste haakjes, dan $\times$
        en $\div$, dan $+$ en $-$);
  \item reken \emph{alleen die} uit en schrijf de hele uitdrukking opnieuw, met
        de uitkomst op zijn plaats;
  \item herhaal tot er één getal overblijft --- één bewerking per regel, geen
        sluiproutes.
\end{enumerate}
\end{method}

\section{Distributiviteit}

\begin{theorem}[Distributiviteit]\label{thm:g7:priorities:distributivity}
Voor alle getallen $k$, $a$ en $b$ geldt:
\[
k \times (a + b) = k \times a + k \times b,
\qquad
k \times (a - b) = k \times a - k \times b .
\]
\end{theorem}

\begin{proof}[Bewijs met een plaatje]
Tel de \omterm{def:g6:measure:area}{oppervlakte} van een rechthoek met breedte $k$ die in twee kleinere
rechthoeken met lengtes $a$ en $b$ verdeeld is: in één geheel is het
$k \times (a+b)$; in twee delen is het $k \times a + k \times b$. Dezelfde
rechthoek, dezelfde \omterm{def:g6:measure:area}{oppervlakte}.
\end{proof}

\begin{omfigure}
\begin{tikzpicture}[scale=0.75]
  \fill[omDef!14] (0,0) rectangle (4.2,1.8);
  \fill[omThm!14] (4.2,0) rectangle (6.4,1.8);
  \draw[thick] (0,0) rectangle (6.4,1.8);
  \draw[thick] (4.2,0) -- (4.2,1.8);
  \node[below, font=\small] at (2.1,0) {$a$};
  \node[below, font=\small] at (5.3,0) {$b$};
  \node[left, font=\small] at (0,0.9) {$k$};
  \node[font=\small] at (2.1,0.9) {$k \times a$};
  \node[font=\small] at (5.3,0.9) {$k \times b$};
\end{tikzpicture}

{\small Distributiviteit als \omterm{def:g6:measure:area}{oppervlaktes}: de grote rechthoek
$k \times (a + b)$ is de \omterm{def:g1:addition:def}{som} van de twee kleinere.}
\end{omfigure}

\begin{example}[Beide richtingen]\label{ex:g7:priorities:distributivity}
\emph{Uitwerken} (van links naar rechts) maakt hoofdrekenen makkelijk:
\[
7 \times 103 = 7 \times (100 + 3) = 700 + 21 = 721,
\qquad
6 \times 98 = 6 \times (100 - 2) = 600 - 12 = 588 .
\]
\emph{Ontbinden} (van rechts naar links) haalt een gemeenschappelijke factor naar
buiten:
\[
23 \times 12 + 23 \times 8 = 23 \times (12 + 8) = 23 \times 20 = 460 .
\]
\end{example}

\begin{example}[Woorden in een berekening omzetten]\label{ex:g7:priorities:words}
``De \omterm{def:g1:addition:def}{som} van $8$ en het \omterm{def:g2:mult:def}{product} van $5$ en $3$'' is $8 + 5 \times 3 = 23$. ``Het
\omterm{def:g2:mult:def}{product} van $5$ en de \omterm{def:g1:addition:def}{som} van $8$ en $3$'' is $5 \times (8 + 3) = 55$: de
woorden ``het \omterm{def:g2:mult:def}{product} van \dots\ en de \omterm{def:g1:addition:def}{som} van\dots'' zetten haakjes om de \omterm{def:g1:addition:def}{som}.
\end{example}

\section{Oefeningen}

\begin{exercise}[$\star$]\label{exo:g7:priorities:1}
Reken uit, één bewerking per regel:
\[
5 + 3 \times 6, \qquad
(5 + 3) \times 6, \qquad
30 - 12 \div 4, \qquad
(30 - 12) \div 4 .
\]
\end{exercise}

\begin{exercise}[$\star$]\label{exo:g7:priorities:2}
Reken uit:
\[
24 - 6 + 2, \qquad
24 \div 6 \times 2, \qquad
24 - 6 \times 2, \qquad
24 \div (6 \times 2).
\]
\end{exercise}

\begin{exercise}[$\star$]\label{exo:g7:priorities:3}
Reken stap voor stap uit:
\[
7 + 3 \times (10 - 6), \qquad
40 - (2 + 3) \times 4, \qquad
60 \div \bigl(2 \times (7 - 4)\bigr).
\]
\end{exercise}

\begin{exercise}[$\star$]\label{exo:g7:priorities:4}
Reken de \omterm{def:g6:fractions:def}{breuken} uit (onzichtbare haakjes!):
\[
\frac{18 + 6}{8}, \qquad
\frac{30}{7 - 2}, \qquad
\frac{5 \times 8}{4 + 6} .
\]
\end{exercise}

\begin{exercise}[$\star$]\label{exo:g7:priorities:5}
Reken met distributiviteit uit het hoofd:
\[
8 \times 102, \qquad
25 \times 99, \qquad
14 \times 12 + 14 \times 8, \qquad
37 \times 45 - 37 \times 35 .
\]
\end{exercise}

\begin{exercise}[$\star$]\label{exo:g7:priorities:6}
Schrijf de berekening op één regel, met haakjes waar dat nodig is, en reken haar
uit: ``het \omterm{ex:g1:subtraction:difference}{verschil} van $50$ en het \omterm{def:g2:mult:def}{product} van $6$ en $7$''; ``het \omterm{def:g2:mult:def}{product} van
de \omterm{def:g1:addition:def}{som} van $4$ en $5$ met het \omterm{ex:g1:subtraction:difference}{verschil} van $10$ en $8$''.
\end{exercise}

\begin{exercise}[$\star$]\label{exo:g7:priorities:7}
Amir kocht $3$ schriften van $2.50$ per stuk en $2$ pennen van $1.20$ per stuk.
Schrijf zijn totaal in één uitdrukking en reken die uit.
\end{exercise}

\begin{exercise}[$\star\star$]\label{exo:g7:priorities:8}
Zet haakjes in $4 + 2 \times 5 - 1$ zodat de uitkomst is: (a) $13$; (b) $29$;
(c) $24$. (Bij een van de drie zijn geen haakjes nodig --- bij welke?)
\end{exercise}

\begin{exercise}[$\star\star$]\label{exo:g7:priorities:9}
Zijn $12 + 4 \times k$ en $16 \times k$ hetzelfde voor elk getal $k$? Test met
$k = 1$ en $k = 2$, en besluit. Welke fout met de voorrangsregels zou van de ene
naar de andere leiden?
\end{exercise}

\begin{exercise}[$\star\star$]\label{exo:g7:priorities:10}
Een bioscoopkaartje kost $9$ euro, maar een groep van $n$ personen betaalt een
vast bedrag van $12$ euro plus $6$ euro per persoon.
\begin{enumerate}
  \item Schrijf een uitdrukking voor de groepsprijs op en reken haar uit voor
        $n = 5$.
  \item Vanaf hoeveel personen is het groepstarief goedkoper dan $n$ losse
        kaartjes? (Probeer een paar waarden van $n$.)
\end{enumerate}
\end{exercise}

\begin{exercise}[$\star\star\star$]\label{exo:g7:priorities:11}
Schrijf met de cijfers $1$, $2$, $3$ en $4$, elk precies één keer, en met de vier
bewerkingen en haakjes naar believen, uitdrukkingen die gelijk zijn aan $24$, aan
$10$ en aan $1$.
\end{exercise}

\section{Opgave: vermenigvuldigen als een Egyptische schrijver}

\begin{problem}[{Weekendopgave --- hoofdrekentrucs uit de distributiviteit, en
het vermenigvuldigen met verdubbelen en halveren dat vierduizend jaar gebruikt
is}]
\label{pb:g7:priorities:1}
Lang voor de rekenmachine --- en voordat de tafels op school werden ingestampt
--- vermenigvuldigden Egyptische schrijvers en Russische boeren elke twee
getallen met alleen \emph{verdubbelen, halveren en optellen}. Hun methode lijkt
magie; ze is in werkelijkheid de distributiviteit
(\cref{thm:g7:priorities:distributivity}) op volle kracht. Deze opgave scherpt
eerst je hoofdrekenen, leert je daarna de oude methode, en opent tot slot de
motorkap om te zien waarom ze precies werkt.

\medskip
\noindent\textbf{Deel I --- De kunst van het niet-rekenen.}
\begin{enumerate}
  \item Reken uit het hoofd uit, door uit te werken zoals in
        \cref{ex:g7:priorities:distributivity}: $7 \times 102$ en
        $9 \times 98$.
  \item Reken uit het hoofd uit, door te ontbinden:
        $17 \times 13 + 17 \times 7$ en $45 \times 101 - 45$.
  \item De ``keer $5$''-truc: om met $5$ te vermenigvuldigen, vermenigvuldig je
        met $10$ en halveer je. Reken $5 \times 87$ zo uit, en leg uit waarom de
        truc mag.
  \item Nog twee trucs om te verantwoorden en te gebruiken: met $25$
        vermenigvuldigen door met $100$ te vermenigvuldigen en door $4$ te delen
        (reken $25 \times 32$ uit); met $99$ vermenigvuldigen door met $100$ te
        vermenigvuldigen en er één keer af te trekken (reken $99 \times 47$
        uit).
  \item Een waarschuwingsschot: distributiviteit verdeelt een factor over een
        \emph{\omterm{def:g1:addition:def}{som}}, nooit over een \emph{\omterm{def:g2:mult:def}{product}}. Vergelijk $8 \times (5 + 3)$
        met $8 \times 5 + 8 \times 3$, en daarna $8 \times (5 \times 3)$ met
        $(8 \times 5) \times (8 \times 3)$: welk paar komt overeen, en welk
        niet?
\end{enumerate}

\medskip
\noindent\textbf{Deel II --- De methode van verdubbelen en halveren.}
Hier is het recept voor $13 \times 24$. Schrijf twee kolommen. Linkerkolom:
begin bij $13$ en \emph{halveer}, waarbij je elke \omterm{def:g3:division:remainder}{rest} laat vallen, tot $1$:
\ $13,\ 6,\ 3,\ 1$. Rechterkolom: begin bij $24$ en \emph{verdubbel} \omterm{def:g3:numbers:evenodd}{even} vaak:
\ $24,\ 48,\ 96,\ 192$. Streep elke rij door waarvan het \emph{linker} getal \omterm{def:g3:numbers:evenodd}{even}
is, en tel de overgebleven rechtergetallen op.
\begin{enumerate}[resume]
  \item Voer het recept uit voor $13 \times 24$: welke rij wordt doorgestreept,
        welke \omterm{def:g1:addition:def}{som} blijft er over, en is die gelijk aan $13 \times 24$?
  \item Vermenigvuldig $21 \times 17$ met de methode (halveer $21$, verdubbel
        $17$) en controleer het antwoord met een gewone vermenigvuldiging.
  \item Vermenigvuldig $16 \times 25$ met de methode. Hoeveel rijen blijven er
        over, en wat is er bijzonder aan een linkerkolom die begint met $1$, $2$,
        $4$, $8$, $16, \dots$?
  \item Nu \emph{waarom} het werkt, stap voor stap. Verantwoord elke gelijkheid
        met distributiviteit:
        \[
        13 \times 24 = (12 + 1) \times 24
        = 12 \times 24 + 24 = 6 \times 48 + 24 .
        \]
        De moraal: links halveren en rechts verdubbelen laat het \omterm{def:g2:mult:def}{product} gelijk
        --- maar als het linkergetal \emph{\omterm{def:g3:numbers:evenodd}{oneven}} is, splitst er één kopie van
        het rechtergetal af die je op de bank moet zetten. Welk getal zet de
        eerste rij van de tabel op de bank?
  \item Ga met de boekhouding verder naar beneden in de tabel:
        $6 \times 48 = 3 \times 96$, dan $3 \times 96 = 1 \times 192 + 96$, en
        dan $1 \times 192 = 192$. Verzamel alles wat je onderweg op de bank hebt
        gezet, en leg uit waarom de bankgetallen precies de rechtergetallen zijn
        van de rijen met een \omterm{def:g3:numbers:evenodd}{oneven} linkergetal --- de rijen die het recept
        bewaart.
\end{enumerate}

\medskip
\noindent\textbf{Deel III --- Onder de motorkap: sommen van verdubbelingen.}
\begin{enumerate}[resume]
  \item In vraag 6 waren de overgebleven rijen $1$, $4$ en $8$ verdubbelingen
        van $24$ waard: de methode schreef in stilte
        $13 \times 24 = (1 + 4 + 8) \times 24$. Ga na dat $1 + 4 + 8 = 13$, en
        zoek welke verdubbelingen ($1$, $2$, $4$, $8$, $16$) de methode voor
        $21$ in vraag 7 bewaarde.
  \item Schrijf $25$ en $42$ als sommen van getallen uit het verdubbelrijtje
        $1, 2, 4, 8, 16, 32$, met elk getal hoogstens één keer.
  \item Leg uit waarom \emph{elk} heel getal zo geschreven kan worden: neem het
        grootste verdubbelgetal dat past, trek het af, en herhaal --- waarom moet
        dit ophouden? Schrijf $100$ als zo'n \omterm{def:g1:addition:def}{som}.
  \item De schrijvers van Egypte deden precies dit, vierduizend jaar geleden: om
        $19 \times 35$ te vermenigvuldigen, maakten ze een verdubbeltabel van
        $35$ ($1 \to 35$, $2 \to 70$, $4 \to 140$, $8 \to 280$, $16 \to 560$),
        kozen de rijen die samen $19$ maken, en telden op. Maak de berekening op
        hun manier.
  \item Finale: moderne computers, die met $0$'en en $1$'en rekenen,
        vermenigvuldigen in wezen als de schrijvers --- met verdubbelingen en
        optellingen. Welke drie optellingen zou een computer maken voor
        $13 \times 24$? (Waarom elk getal in verdubbelingen uiteenvalt, en wat
        $0$'en en $1$'en daarmee te maken hebben, is de weekendopgave van het
        hoofdstuk over machten: \cref{pb:g8:powers:1}.)
\end{enumerate}
\end{problem}
